% =============================================================================
%  capitulos/cap7_conclusao.tex
%  CAPÍTULO 7 — CONCLUSÃO
%
%  Objetivo: amarrar o trabalho. Olhar para trás (o que foi feito e
%  se os objetivos foram atingidos) e para frente (o que pode melhorar).
%
%  REGRAS IMPORTANTES:
%    ✗ Não introduza conteúdo novo aqui (sem novos conceitos ou dados)
%    ✗ Não repita trechos inteiros de capítulos anteriores
%    ✓ Retome os objetivos específicos um a um e avalie cada um
%    ✓ Seja honesto sobre o que não foi atingido e por quê
% =============================================================================

\chapter{Conclusão}
\label{cap:conclusao}


% -----------------------------------------------------------------------------
%  7.1 CONSIDERAÇÕES FINAIS
%  Retome o problema inicial e avalie se os objetivos foram alcançados.
%  Para cada objetivo específico do Cap. 1, diga: foi atingido? Totalmente?
%  Parcialmente? Justifique brevemente.
% -----------------------------------------------------------------------------
\section{Considerações Finais}
\label{sec:consideracoes-finais}

TODO: Comece retomando o problema apresentado na introdução e avalie
se o objetivo geral foi alcançado. Em seguida, percorra os objetivos
específicos:

\begin{itemize}
  \item TODO: O objetivo ``[copie o objetivo específico 1]'' foi
        \textbf{atingido}, conforme demonstrado em [Seção X / Figura Y];
  \item TODO: O objetivo ``[copie o objetivo específico 2]'' foi
        \textbf{parcialmente atingido}, pois [justificativa];
  \item TODO: Continue para cada objetivo específico listado no Cap. 1.
\end{itemize}

TODO: Finalize com um parágrafo destacando as principais contribuições
do trabalho — o que ele entrega de valor para o usuário ou para a área.


% -----------------------------------------------------------------------------
%  7.2 DIFICULDADES ENCONTRADAS
%  Relate os obstáculos reais enfrentados durante o desenvolvimento.
%  Seja honesto. Isso demonstra maturidade e reflexão crítica.
%  Ex.: dificuldades técnicas, limitações de tempo, falta de dados,
%       mudanças de escopo, problemas de integração com APIs externas...
% -----------------------------------------------------------------------------
\section{Dificuldades Encontradas}
\label{sec:dificuldades}

TODO: Descreva os principais obstáculos técnicos, metodológicos ou de
contexto enfrentados ao longo do projeto e como foram superados (ou
contornados). Exemplos:

\begin{itemize}
  \item \textbf{TODO: Dificuldade técnica}: ex.: integração com a API X
        apresentou limitações de rate limit, resolvidas com [abordagem];
  \item \textbf{TODO: Limitação de tempo}: ex.: algumas funcionalidades
        planejadas precisaram ser descartadas para viabilizar a entrega;
  \item \textbf{TODO: Outras dificuldades}: adicione conforme o seu contexto.
\end{itemize}


% -----------------------------------------------------------------------------
%  7.3 TRABALHOS FUTUROS
%  O que ficou de fora do escopo atual? Como o sistema pode evoluir?
%  O que a metodologia poderia ter feito melhor?
%  Esta seção demonstra visão além do entregável imediato.
% -----------------------------------------------------------------------------
\section{Trabalhos Futuros}
\label{sec:trabalhos-futuros}

A partir da experiência acumulada no desenvolvimento deste trabalho,
identificam-se as seguintes oportunidades de evolução:

\begin{itemize}
  \item TODO: Implementar [funcionalidade planejada mas não executada]
        para ampliar a cobertura de [RF ou objetivo];
  \item TODO: Realizar testes de usabilidade com uma amostra maior e
        mais diversificada de usuários reais;
  \item TODO: Migrar a infraestrutura para um ambiente de nuvem (ex.:
        AWS, GCP, Azure) para garantir maior escalabilidade e
        disponibilidade;
  \item TODO: Integrar o sistema com [serviço ou API externa] para
        enriquecer as funcionalidades de [área];
  \item TODO: Adicionar suporte a [plataforma, idioma, acessibilidade...];
  \item TODO: Outros pontos de melhoria identificados ao longo do projeto.
\end{itemize}
